% ============================================================
% Drop-in replacement for Section 3 (Theoretical Framework) of main.tex.
% Adds sanitary visits / regulator information to the Kang--Silveira CWS model:
%   - continuous health-risk type r, unobserved by the regulator (CWS knows its own r)
%   - costly sanitary visits Q(a) reveal r; visits are triggered by violations
%   - regulator is risk neutral when it does not know the health type
%     (no-visit penalty is built on the average risk)
% Usage: \input this file in place of the current Section 3, or copy its body in.
% Requires only amsmath + amssymb, already loaded in main.tex (\mathbb needs amssymb).
% Nothing in main.tex is modified by creating this file.
% ============================================================

\section{Theoretical Framework}

We model the effect of coal mining pollution on a CWS's decision to comply with the Safe Drinking Water Act under regulatory enforcement by adapting a static model of environmental compliance from \cite{kang2021understanding}, which is itself an adaptation of \cite{mookherjee1994marginal}. We extend that framework in one direction that is central to our setting: the regulator does not observe the health harm of a violation and uses costly sanitary visits to assess it. Because a monitoring and reporting violation is the \emph{absence} of a reported test, the regulator cannot read the contamination behind it; the sanitary visit is the instrument it uses to learn how harmful a given violation is. Embedding this single feature lets the model speak jointly to the three margins we document empirically: violations, sanitary visits, and the type of enforcement that follows.

\subsection{The CWS compliance problem}

CWSs choose a negligence level $a \in [0,\bar{a}]$, which measures how much they avoid the requirements of drinking water quality regulation, from maximum negligence $\bar{a}$ to full compliance $0$. Failing to conduct or submit tests, or not installing pollution abatement technology, are examples of negligence. Regulators do not observe $a$. CWSs face a compliance cost $\theta$, a random variable drawn from $\Theta$ with strictly increasing, continuously differentiable distribution $F(\cdot)$ on support $(0,\bar{\theta})$. CWSs know their draw of $\theta$; regulators know only the distribution. A compliance cost reflects both a CWS's ability to manage a water system and the cost of water treatment.

A CWS receives benefit $\theta b(a)$ from compliance cost $\theta$ and negligence $a$, so the opportunity cost of choosing $a$ relative to maximum negligence is $\theta[b(\bar{a})-b(a)]$. The number of violations is a random variable $K$ that is Poisson with mean $a$,
\[
\Pr(K=k\mid a) = e^{-a}\,\frac{a^k}{k!}.
\]

\subsection{Regulator information and sanitary visits}

We add a second, continuous type. Each CWS has a \emph{health-risk type} $r$, the social and health harm of a violation at that system, drawn from a distribution $G(\cdot)$ with density $g(\cdot)$ on support $(0,\bar{r})$ and mean $\mu_r \equiv \mathbb{E}[r]$. A high-$r$ system is one where a failure to monitor is more dangerous, for example because its source water is more contaminated, its treatment margin thinner, or its served population larger. Following the costly-state-verification logic, the CWS knows its own $r$; the regulator does not, unless it conducts a sanitary visit.

A sanitary visit reveals $r$. Visits are costly and are triggered by violations: the ex ante probability that a CWS with negligence $a$ is visited is $Q(a)$, with $Q'(a)>0$. This captures that the regulator, facing resource constraints, prioritizes inspections at systems that have accumulated monitoring and reporting violations.

The expected penalty depends on the regulator's information state. If the CWS is \emph{not} visited, the regulator cannot condition on $r$ and, being risk neutral, sets the penalty at the average risk $\mu_r$. Write this no-visit expected penalty as
\[
e_N(a) = e^{-a}\sum_{k=0}^{\infty} \frac{\epsilon_N(k)}{k!}\,a^k ,
\]
where $\epsilon_N(\cdot)$ is the schedule the regulator applies when it knows only the distribution of $r$. If the CWS \emph{is} visited and risk $r$ is revealed, the regulator conditions on it, yielding $e_V(a,r)$ with $\partial e_V/\partial r \ge 0$: a higher revealed risk warrants a stricter penalty. We normalize $e_V(a,\mu_r) = e_N(a)$, so that learning ``average'' risk leaves the penalty unchanged.

Define the \emph{risk gap}
\[
D(a,r) \equiv e_V(a,r) - e_N(a), \qquad \operatorname{sign} D(a,r) = \operatorname{sign}(r-\mu_r), \qquad D_r \equiv \tfrac{\partial D}{\partial r} > 0 .
\]
A CWS revealed to be high risk ($r>\mu_r$) is penalized more than the average ($D>0$); one revealed to be low risk ($r<\mu_r$) is penalized less ($D<0$). The CWS, knowing its own $r$, faces expected penalty
\begin{equation}\label{eq:expected_penalty_visit}
P(a,r) = [1-Q(a)]\,e_N(a) + Q(a)\,e_V(a,r) = e_N(a) + Q(a)\,D(a,r).
\end{equation}

The expected payoff to a CWS of type $(\theta,r)$ from negligence $a$ is
\begin{equation}\label{eq:cws_payoff_visit}
-\theta[b(\bar{a})-b(a)] - P(a,r).
\end{equation}
A negligence schedule $a(\cdot)$ is implemented by the penalty environment if $a(\theta,r)$ maximizes \eqref{eq:cws_payoff_visit}. Combining the first order condition with the implemented schedule gives
\begin{equation}\label{eq:foc_visit}
\theta\, b'(a) = M(a,r), \qquad M(a,r) \equiv e_N'(a) + Q'(a)\,D(a,r) + Q(a)\,D_a(a,r),
\end{equation}
whenever $a(\theta,r)>0$, where $M(a,r)$ is the \emph{marginal expected penalty}. The new term $Q'(a)\,D(a,r)$ is the deterrence created by the visit process: raising negligence raises the chance of triggering a visit, and a visit moves the penalty by $D$. For a high-risk CWS this term \emph{adds} to the marginal penalty (a visit would raise the penalty); for a low-risk CWS it \emph{subtracts} (a visit would lower it). We assume $b$ is strictly increasing with $b''<0$, that $M$ is increasing in $a$, and that the first order condition characterizes the maximum, so the second order condition
\[
\theta\, b''(a) - M_a(a,r) < 0, \qquad\text{equivalently}\qquad \Phi \equiv M_a(a,r) - \theta\, b''(a) > 0,
\]
holds. We use $\Phi>0$ throughout. When $Q\equiv 0$ the model collapses to \cite{kang2021understanding} with $M=e'$.

\subsection{Revised comparative statics}

\noindent\textbf{Proposition 1.} \emph{CWSs choose higher negligence as compliance costs rise, $\partial a/\partial\theta > 0$, and the strength of the response depends on the CWS's health-risk type.}

Differentiating \eqref{eq:foc_visit} in $\theta$,
\[
\frac{\partial a}{\partial\theta} = -\frac{b'(a)}{\theta\, b''(a) - M_a(a,r)} = \frac{b'(a)}{\Phi} > 0 .
\]
The sign is the same as in the model without visits, but $\Phi$ now contains the visit terms and depends on $r$ through $M_a$. The negligence response to compliance cost is steeper for low-risk systems, where visits do not deter, and flatter for high-risk systems. The original result is the special case $Q\equiv 0$.

\medskip
\noindent\textbf{Proposition 2.} \emph{Suppose mining raises contamination so that compliance cost is nondecreasing in upstream mining, $\theta(m)$ with $\theta'(m)\ge 0$, and also raises the health stakes of a violation, $r(m)$ with $r'(m)\ge 0$. Then}
\[
\frac{\partial a}{\partial m} = \frac{\partial a}{\partial\theta}\,\theta'(m) + \frac{\partial a}{\partial r}\,r'(m) = \frac{1}{\Phi}\Big[\underbrace{b'(a)\,\theta'(m)}_{\text{cost channel }(+)} - \underbrace{M_r\,r'(m)}_{\text{deterrence channel }(-)}\Big],
\]
\emph{which is positive if and only if the cost channel dominates.}

The own-risk response is obtained by differentiating \eqref{eq:foc_visit} in $r$,
\[
\frac{\partial a}{\partial r} = \frac{M_r}{\theta\, b''(a) - M_a(a,r)} = -\frac{M_r}{\Phi} < 0, \qquad M_r = Q'(a)\,D_r + Q(a)\,D_{ar} > 0 ,
\]
so a higher own health risk \emph{lowers} negligence: a high-risk CWS, knowing a visit would expose it to a stricter penalty and that visits are triggered by its own violations, complies more. Mining therefore acts through two opposing forces. Through compliance cost, $\theta'(m)$ raises the marginal benefit of negligence and pushes $a$ up; through health stakes, $r'(m)$ makes visits more consequential and pulls $a$ down. The original Proposition~2 is the special case in which $r$ is fixed. Our empirical finding $\partial a/\partial m > 0$ then carries content it did not before: it reveals that the cost channel dominates the regulator's visit-based deterrence---the inspection response is too weak to offset the rising benefit of non-compliance.

\medskip
\noindent\textbf{Proposition 3.} \emph{Stiffer penalties reduce negligence; the effect of more sanitary visits depends on the CWS's health-risk type.}

Scaling all penalties by $\lambda>1$, so $P \to \lambda P$ and the first order condition is $\theta b'(a) = \lambda M(a,r)$,
\[
\frac{\partial a}{\partial\lambda} = \frac{M(a,r)}{\theta\, b''(a) - \lambda M_a(a,r)} = -\frac{M(a,r)}{\lambda M_a - \theta b''} < 0
\]
for systems with $r \ge \mu_r$ (where $M>0$), reproducing the original result. A separate lever is the \emph{intensity} of inspections. Scaling the visit probability $Q \to \mu Q$,
\[
\frac{\partial a}{\partial\mu} = -\frac{M_Q}{\Phi}, \qquad M_Q \equiv Q'(a)\,D(a,r) + Q(a)\,D_a(a,r), \qquad \operatorname{sign}\frac{\partial a}{\partial\mu} = -\operatorname{sign}(r-\mu_r) .
\]
More sanitary visits raise compliance among high-risk systems ($r>\mu_r$), which a visit would expose, but \emph{lower} it among low-risk systems, which a visit would exonerate. Blanket increases in inspection therefore improve compliance only where health risk is genuinely high.

\subsection{Regulator information and sanitary visits: testable predictions}

The visit process generates a set of comparative statics on objects we observe directly---violations, visits, and enforcement type---that the model without visits is silent on.

\medskip
\noindent\textbf{Proposition A (visits rise in violations).} The inspection technology satisfies $Q'(a) > 0$: a CWS that accumulates more violations is more likely to be visited.

\medskip
\noindent\textbf{Proposition B (mining raises visits).} Composing Proposition~A with Proposition~2,
\[
\frac{\partial Q}{\partial m} = Q'(a)\,\frac{\partial a}{\partial m} > 0 ,
\]
so upstream mining raises the probability of a sanitary visit through its effect on violations.

\medskip
\noindent\textbf{Proposition C (mining lowers formal enforcement).} Formal enforcement is the strict, high-$r$ branch of $e_V$. Because the risk a visit \emph{assesses} does not rise with mining in our sample---measured concentrations remain below the maximum contaminant level---and because mining shifts the violation mix toward monitoring and reporting violations, which are handled informally, the probability of formal enforcement falls with mining,
\[
\frac{\partial\,\Pr(\text{formal}\mid\text{visit})}{\partial m} < 0 .
\]

\medskip
\noindent\textbf{Proposition D (visits route to informal enforcement).} A visit moves the graduated, informal branch of $e_V$ rather than the rare formal branch,
\[
\frac{\partial\,\Pr(\text{informal}\mid\text{visit})}{\partial(\text{visit})} > 0, \qquad \frac{\partial\,\Pr(\text{formal}\mid\text{visit})}{\partial(\text{visit})} \approx 0 .
\]

\medskip
\noindent\textbf{Proposition E (heterogeneous compliance in own risk).} From Proposition~2, $\partial a/\partial r < 0$: systems with higher health stakes comply more, conditional on compliance cost. This is the one prediction not yet tested directly; it can be examined by interacting the violation response with proxies for health stakes such as population served, surface versus groundwater source, or baseline contamination.

\medskip
Taken together, Propositions~B and~C expose the wedge that organizes our results. Let $r(m)$ be the \emph{true} health harm of a violation, rising with mining, and let $\hat{r}(m)$ be the risk a sanitary visit actually \emph{recovers}. Visits rise with mining (Proposition~B) while formal enforcement falls (Proposition~C), so $\hat{r}(m)$ is flat or falling even as $r(m)$ rises. The reason is that a visit samples currently available, tested water, which lies below the maximum contaminant level, precisely when the worst contamination is concentrated in the unmonitored periods a monitoring and reporting violation conceals. The sanitary visit, the regulator's information technology for $r$, systematically under-recovers harm in exactly the strategic states. This is also the model's point of departure from \cite{kang2021understanding}: there the harm $h(a)$ is a primitive the regulator observes for free, whereas here it must be acquired through costly visits, and the acquisition is biased in the non-monitoring states.

\subsection{Mapping propositions to the empirical results}

Each proposition is tested by a specific table. Proposition~1 (negligence rises with compliance cost) is supported jointly by the concentration regressions in Table~\ref{tab:6yr_huc02fe_inorg_ravalli_2005}, which show that mining raises compliance cost by raising contaminant concentrations, and by the violation regressions in Table~\ref{2sls_dwnstrm_minevio_mr}, which show the resulting rise in monitoring and reporting violations. Proposition~2 (mining raises negligence, cost channel dominating the visit-based deterrence) is tested in Table~\ref{2sls_dwnstrm_minevio_mr}, where an additional upstream mine raises monitoring and reporting violations by roughly 14 to 20 days per year, alongside the null effect on maximum contaminant level violations in Table~\ref{2sls_dwnstrm_minevio_mcl}. Proposition~3 (lower penalties raise negligence) is supported by the fall in formal enforcement with mining in Table~\ref{tab:h3_inf_formal_d12}. Proposition~A ($Q'(a)>0$) is tested by the regression of sanitary visits on monitoring and reporting violations in Table~\ref{tab:mr_to_visit}. Proposition~B (mining raises visits) is the 14 percentage point effect in Table~\ref{tab:h2_snsv_d12}. Proposition~C (mining lowers formal enforcement) is the 5.6 percentage point fall in Table~\ref{tab:h3_inf_formal_d12}. Proposition~D (visits route to informal enforcement) is tested by the regression of enforcement type on sanitary visits in Table~\ref{tab:visit_to_enf}, which shows visits raising informal but not formal enforcement. Proposition~E ($\partial a/\partial r < 0$) is not yet tested and is the natural next specification, interacting the violation response with proxies for health stakes.
